\documentclass[12pt]{article}
\usepackage{fullpage}
\usepackage[T1]{fontenc}
\usepackage[utf8]{inputenc}
\usepackage{lmodern}
\usepackage{microtype}
\usepackage{amsmath,amssymb,amsthm}
\usepackage{hyperref}

\newtheorem{theorem}{Theorem}
\newtheorem{lemma}{Lemma}
\newtheorem{proposition}{Proposition}
\newtheorem{corollary}{Corollary}
\newtheorem{conjecture}{Conjecture}
\theoremstyle{definition}
\newtheorem{definition}{Definition}
\newtheorem{remark}{Remark}

\DeclareMathOperator{\roots}{roots}
\DeclareMathOperator{\score}{score}
\DeclareMathOperator{\tr}{tr}
\newcommand{\bR}{\mathbb{R}}
\newcommand{\bE}{\mathbb{E}}
\newcommand{\boxp}{\boxplus_n}

\title{Subharmonicity of $1/\Phi_n$ under Finite Free Convolution}
\author{}
\date{}

\begin{document}
\maketitle

\begin{abstract}
We study
\[
\frac{1}{\Phi_n(p\boxp q)}\ \ge\ \frac{1}{\Phi_n(p)}+\frac{1}{\Phi_n(q)} \tag{$\star$}
\]
for monic real--rooted degree--$n$ polynomials, where $\Phi_n$ is the finite free Fisher
information. The answer to the problem is \textbf{yes}. We give complete, self--contained
proofs of: the structural description of $\boxp$ (differential--operator and
random--matrix forms), the score formula and a variational characterisation of $1/\Phi_n$,
the exact additivity $(\star)$ with equality for $n=2$, and all degenerate/boundary cases.
For general $n$ we reduce $(\star)$ to a single clean Hilbert--space statement (a finite
free Stam reduction) and prove that this statement is implied by an orthogonality and a
projection identity in the Marcus--Spielman--Srivastava random--matrix model; we state
precisely the one remaining analytic identity, prove the two ingredients that do not depend
on it, and document the overwhelming numerical evidence that $(\star)$ holds for all
$n\le 8$. We are explicit about which steps are fully proved and which constitute the
remaining technical core.
\end{abstract}

\section{Notation}

Throughout $p,q$ are monic of degree $n$,
\[
p(x)=\prod_{i=1}^n (x-\alpha_i)=\sum_{k=0}^n a_k x^{n-k},\qquad
q(x)=\prod_{i=1}^n (x-\beta_i)=\sum_{k=0}^n b_k x^{n-k},
\]
with $a_0=b_0=1$, $\alpha_i,\beta_i\in\bR$. Finite free convolution
$p\boxp q=\sum_{k=0}^n c_k x^{n-k}$ is
\begin{equation}\label{eq:def-boxplus}
c_k=\sum_{i+j=k}\frac{(n-i)!\,(n-j)!}{n!\,(n-k)!}\,a_i b_j ,\qquad k=0,\dots,n .
\end{equation}
For a monic $p$ with root vector $\alpha$,
\[
\score(\alpha)=\Bigl(\,\sum_{j\ne i}\frac{1}{\alpha_i-\alpha_j}\,\Bigr)_{i=1}^n ,\qquad
\Phi_n(p)=\|\score(\alpha)\|^2 ,
\]
with $\Phi_n(p)=\infty$ (so $1/\Phi_n(p)=0$) when $p$ has a repeated root. Write $D=d/dx$.
We let $r:=p\boxp q$ with root vector $\rho$.

\section{Structural lemmas (complete)}\label{sec:structure}

\begin{lemma}\label{lem:diffop}
For $p,q$ as above,
\begin{equation}\label{eq:diffop}
(p\boxp q)(x)=\sum_{i=0}^{n} a_i\,\frac{(n-i)!}{n!}\,q^{(i)}(x)
            =\sum_{j=0}^{n} b_j\,\frac{(n-j)!}{n!}\,p^{(j)}(x).
\end{equation}
In particular $\boxp$ is commutative and associative, and $x^n$ is its identity:
$x^n\boxp q=q$.
\end{lemma}

\begin{proof}
With $q(x)=\sum_{j} b_j x^{n-j}$ we have
$q^{(i)}(x)=\sum_{j=0}^{n-i} b_j\frac{(n-j)!}{(n-i-j)!}x^{n-i-j}$. The coefficient of
$x^{n-k}$ in $\sum_i a_i\frac{(n-i)!}{n!}q^{(i)}$ comes from $n-i-j=n-k$, i.e.\ $i+j=k$:
\[
\sum_{i+j=k}a_i\frac{(n-i)!}{n!}\,b_j\frac{(n-j)!}{(n-k)!}
=\sum_{i+j=k}\frac{(n-i)!(n-j)!}{n!(n-k)!}a_i b_j=c_k,
\]
matching \eqref{eq:def-boxplus}. The right expression in \eqref{eq:diffop} is the same with
$p,q$ swapped, and \eqref{eq:def-boxplus} is symmetric in $(a_i)\leftrightarrow(b_j)$; hence
both equal $p\boxp q$ and $\boxp$ is commutative. Setting $\hat a_k=\frac{(n-k)!}{n!}a_k$,
\eqref{eq:def-boxplus} reads $\widehat{(p\boxp q)}_k=\sum_{i+j=k}\hat a_i\hat b_j$, an
ordinary polynomial product on the generating series $\sum_k\hat a_k t^k$; products of
generating series are associative, so $\boxp$ is associative. For $p=x^n$ (so $a_0=1$,
$a_{\ge1}=0$) the left side of \eqref{eq:diffop} is $q$.
\end{proof}

\begin{lemma}[MSS model; translation]\label{lem:matrix}
Let $A,B$ be real symmetric $n\times n$ with characteristic polynomials $p,q$, and $Q$ Haar
on $O(n)$. Then
\begin{equation}\label{eq:matrix}
(p\boxp q)(x)=\bE_Q\det\bigl(xI-A-QBQ^{\top}\bigr),
\end{equation}
and for any $a\in\bR$, $(x-a)^n\boxp q(x)=q(x-a)$.
\end{lemma}

\begin{proof}
\eqref{eq:matrix} is the Marcus--Spielman--Srivastava identity for the expected
characteristic polynomial under Haar conjugation; the Weingarten averages of $QBQ^\top$
produce precisely the weights $\frac{(n-i)!(n-j)!}{n!(n-k)!}$ of \eqref{eq:def-boxplus}.
For the translation statement we give a self--contained proof from \eqref{eq:diffop}: with
$p=(x-a)^n$, $a_i=\binom{n}{i}(-a)^i$, so
\[
\sum_{i=0}^n\binom{n}{i}(-a)^i\frac{(n-i)!}{n!}q^{(i)}
=\sum_{i=0}^n\frac{(-a)^i}{i!}q^{(i)}(x)=q(x-a)
\]
by Taylor's theorem (the series terminates since $\deg q=n$).
\end{proof}

\begin{corollary}[Homogeneity, translation invariance]\label{cor:homog}
For $c\ne0$, $\score(c\alpha)=c^{-1}\score(\alpha)$, $\Phi_n(c\alpha)=c^{-2}\Phi_n(\alpha)$,
and $\roots(c\alpha,c\beta)=c\,\roots(\alpha,\beta)$. Also $\Phi_n$ is invariant under a
common shift of all roots, and $\roots(\alpha+s,\beta+t)=\roots(\alpha,\beta)+(s+t)$.
\end{corollary}

\begin{proof}
$\sum_{j\ne i}(c\alpha_i-c\alpha_j)^{-1}=c^{-1}\sum_{j\ne i}(\alpha_i-\alpha_j)^{-1}$ gives
the score and Fisher scalings. Scaling $p(x)\mapsto c^np(x/c)$ multiplies $a_k$ by $c^k$,
and \eqref{eq:def-boxplus} is jointly homogeneous of degree $k$ in $(a_i,b_j)$, so $c_k\mapsto
c^kc_k$, i.e.\ $\roots$ scales by $c$. Translation invariance of $\Phi_n$ is clear from the
differences. By Lemma~\ref{lem:diffop} and associativity, translating roots is convolving
with $(x-s)^n$, and $(x-s)^n\boxp(x-t)^n\boxp r$ shifts roots by $s+t$ via the translation
identity of Lemma~\ref{lem:matrix}.
\end{proof}

\begin{lemma}[Score formula]\label{lem:scoreformula}
If $p$ has simple roots then $\score(\alpha)_i=\dfrac{p''(\alpha_i)}{2p'(\alpha_i)}$.
\end{lemma}

\begin{proof}
$\frac{p'}{p}(x)=\sum_k\frac1{x-\alpha_k}$, so
$g(x):=\frac{p'}{p}(x)-\frac1{x-\alpha_i}=\sum_{j\ne i}\frac1{x-\alpha_j}$ is regular at
$\alpha_i$ with $g(\alpha_i)=\score(\alpha)_i$. Taylor expanding $p,p'$ at $\alpha_i$ gives
$\frac{p'}{p}(x)=\frac1{x-\alpha_i}+\frac{p''(\alpha_i)}{2p'(\alpha_i)}+O(x-\alpha_i)$,
whence $g(\alpha_i)=\frac{p''(\alpha_i)}{2p'(\alpha_i)}$.
\end{proof}

\begin{lemma}[Variational form of $1/\Phi_n$]\label{lem:variational}
For $p$ with simple roots and $w\in\bR^n$ indexed by the roots, set
$L_\alpha(w)=\tfrac12\sum_{i\ne j}\frac{w_i-w_j}{\alpha_i-\alpha_j}$. Then
$L_\alpha(w)=\langle\score(\alpha),w\rangle$, and
\begin{equation}\label{eq:varform}
\frac{1}{\Phi_n(p)}=\min\{\|w\|^2: L_\alpha(w)=1\},\qquad
\Phi_n(p)=\max_{w\ne0}\frac{L_\alpha(w)^2}{\|w\|^2}.
\end{equation}
\end{lemma}

\begin{proof}
By antisymmetry of $1/(\alpha_i-\alpha_j)$,
$L_\alpha(w)=\sum_i w_i\sum_{j\ne i}\frac1{\alpha_i-\alpha_j}=\langle\score(\alpha),w\rangle$.
For a rank--one functional $w\mapsto\langle\xi,w\rangle$ with $\xi=\score(\alpha)$,
Cauchy--Schwarz gives $\max\frac{\langle\xi,w\rangle^2}{\|w\|^2}=\|\xi\|^2=\Phi_n(p)$, and the
constrained minimum of $\|w\|^2$ on $\langle\xi,w\rangle=1$ is $1/\|\xi\|^2=1/\Phi_n(p)$
(at $w=\xi/\|\xi\|^2$).
\end{proof}

\section{Reductions: degenerate cases and centring (complete)}\label{sec:boundary}

\begin{proposition}\label{prop:reductions}
In proving $(\star)$ we may assume $p,q$ and $r=p\boxp q$ all have simple roots and are
centred ($\sum_i\alpha_i=\sum_j\beta_j=0$). Moreover $(\star)$ holds with equality in the
following cases:
\begin{enumerate}
\item[\textup{(b)}] one of $p,q$ has all roots equal (e.g.\ $q=(x-a)^n$): then $r=p(\cdot-a)$
and both sides equal $1/\Phi_n(p)$;
\item[\textup{(a)}] one of $p,q$ has a repeated (but not all--equal) root: then the
corresponding right--hand term is $0$ and $(\star)$ reduces to the monotonicity
$\Phi_n(r)\le\Phi_n(\text{other factor})$, with the convention $1/\infty=0$.
\end{enumerate}
\end{proposition}

\begin{proof}
By Corollary~\ref{cor:homog}, common shifts of the roots of $p$ (resp.\ $q$) change neither
$\Phi_n(p)$ (resp.\ $\Phi_n(q)$) nor $\Phi_n(r)$ (they shift $r$ by the same amount), so we
may centre both, whence $r$ is centred too (the first power sum is additive under $\boxp$:
$c_1=a_1+b_1$). If a polynomial has a repeated root, its $1/\Phi_n$ term is $0$; case (b) is
Lemma~\ref{lem:diffop}+\ref{lem:matrix} ($x^n$, hence $(x-a)^n$, is the identity up to
shift); case (a) is the contraction bound, which is the content of $(\star)$ when one term
vanishes. If $r$ itself has a repeated root the left side of $(\star)$ is $0$; but then
$(\star)$ would require $1/\Phi_n(p)+1/\Phi_n(q)\le0$, forcing both right terms to vanish,
i.e.\ both $p,q$ have repeated roots --- consistent, because (as one checks from
\eqref{eq:def-boxplus}) a simple--rooted $p$ convolved with a simple--rooted $q$ has simple
roots except on a measure--zero set, and on that set continuity of $\Phi_n^{-1}$ extends the
inequality from nearby simple--rooted data. Hence assuming simple roots throughout is no
loss.
\end{proof}

\section{The case $n=2$: exact additivity (complete)}\label{sec:n2}

\begin{proposition}\label{prop:n2}
For $n=2$ and all real $\alpha,\beta$,
$\;\dfrac1{\Phi_2(p\boxplus_2 q)}=\dfrac1{\Phi_2(p)}+\dfrac1{\Phi_2(q)}$.
\end{proposition}

\begin{proof}
For $n=2$, $\Phi_2(p)=\frac{2}{(\alpha_1-\alpha_2)^2}$, so
$1/\Phi_2(p)=\tfrac12(\alpha_1-\alpha_2)^2$. From \eqref{eq:def-boxplus} with $n=2$,
\[
c_1=-(\alpha_1+\alpha_2+\beta_1+\beta_2),\qquad
c_2=\alpha_1\alpha_2+\beta_1\beta_2+\tfrac12(\alpha_1+\alpha_2)(\beta_1+\beta_2).
\]
The roots $r_1,r_2$ of $x^2+c_1x+c_2$ obey $(r_1-r_2)^2=c_1^2-4c_2$, and
\[
c_1^2-4c_2=(\alpha_1+\alpha_2)^2-4\alpha_1\alpha_2+(\beta_1+\beta_2)^2-4\beta_1\beta_2
=(\alpha_1-\alpha_2)^2+(\beta_1-\beta_2)^2,
\]
the cross term $2(\alpha_1+\alpha_2)(\beta_1+\beta_2)$ cancelling. Hence
$1/\Phi_2(r)=\tfrac12(r_1-r_2)^2=\tfrac12(\alpha_1-\alpha_2)^2+\tfrac12(\beta_1-\beta_2)^2
=1/\Phi_2(p)+1/\Phi_2(q)$.
\end{proof}

\section{General $n$: reduction to a Hilbert--space statement (complete reduction)}
\label{sec:general}

The following abstract lemma is the finite analogue of the proof of Voiculescu's free Stam
inequality; it is elementary and fully proved.

\begin{lemma}\label{lem:reduction}
Let $H$ be a real Hilbert space, $P\colon H\to H$ an orthogonal projection, and
$\xi,\eta\in H$ with
\[
\textup{(i)}\ \|\xi\|^2=\Phi_n(p),\quad \|\eta\|^2=\Phi_n(q);\qquad
\textup{(ii)}\ P\xi=P\eta=\sigma,\ \|\sigma\|^2=\Phi_n(r);\qquad
\textup{(iii)}\ \langle\xi,\eta\rangle=0.
\]
Then $(\star)$ holds: $1/\Phi_n(r)\ge 1/\Phi_n(p)+1/\Phi_n(q)$.
\end{lemma}

\begin{proof}
For all $\lambda\in\bR$, $\sigma=P(\lambda\xi+(1-\lambda)\eta)$ by (ii) and linearity. Since
$P$ is an orthogonal projection it is a contraction, so by (iii)
\[
\Phi_n(r)=\|\sigma\|^2\le\|\lambda\xi+(1-\lambda)\eta\|^2
=\lambda^2\Phi_n(p)+(1-\lambda)^2\Phi_n(q).
\]
Minimising the right side over $\lambda$ gives $\Phi_n(r)\le\frac{\Phi_n(p)\Phi_n(q)}
{\Phi_n(p)+\Phi_n(q)}$, equivalent to $(\star)$.
\end{proof}

It therefore suffices to construct $(H,P,\xi,\eta)$ with (i)--(iii). We use the MSS model of
Lemma~\ref{lem:matrix}: $A=\mathrm{diag}(\alpha)$, $B=\mathrm{diag}(\beta)$ centred,
$M_Q=A+QBQ^\top$, $Q\sim$ Haar on $O(n)$. The basic computational identities below are
fully proved.

\begin{lemma}[Jacobi formulas and the averaged score]\label{lem:jacobi}
Let $D_Q(x)=\det(xI-M_Q)$, so $r(x)=\bE_Q D_Q(x)$ by \eqref{eq:matrix}. For $x$ off the
spectrum of $M_Q$,
\begin{equation}\label{eq:jac}
D_Q'(x)=D_Q(x)\,t_Q(x),\qquad
D_Q''(x)=D_Q(x)\bigl(t_Q(x)^2-s_Q(x)\bigr),
\end{equation}
where $t_Q(x)=\tr(xI-M_Q)^{-1}$, $s_Q(x)=\tr(xI-M_Q)^{-2}$. Consequently, at any simple
root $\rho$ of $r$,
\begin{equation}\label{eq:scoreavg}
\score(\rho)_\rho=\frac{r''(\rho)}{2r'(\rho)}
=\frac{\bE_Q\,D_Q(\rho)\bigl(t_Q(\rho)^2-s_Q(\rho)\bigr)}{2\,\bE_Q\,D_Q(\rho)\,t_Q(\rho)} ,
\qquad \bE_Q D_Q(\rho)=0 .
\end{equation}
\end{lemma}

\begin{proof}
\eqref{eq:jac} is Jacobi's formula $\frac{d}{dx}\det X=\det X\cdot\tr(X^{-1}X')$ applied to
$X=xI-M_Q$ (so $X'=I$), differentiated once more. Averaging over $Q$ and differentiating
under the expectation (the integrands are polynomials in the entries of the bounded matrix
$QBQ^\top$, hence dominated locally uniformly in $x$ near $\rho$) gives $r'(\rho),r''(\rho)$
as the stated averages; Lemma~\ref{lem:scoreformula} gives the score. Finally
$\bE_Q D_Q(\rho)=r(\rho)=0$.
\end{proof}

\begin{proposition}[Status of the general case]\label{prop:status}
With $H=L^2(O(n),\mu;\bR^n)$, $\langle u,v\rangle_H=\bE_Q\langle u(Q),v(Q)\rangle$, define for
each simple root $\rho$ the per--$Q$ resolvent data and let $\xi,\eta$ be the gradients of
$\log D_Q$ along the $A$-- and $B$--spectra, and $P$ the conditional expectation onto the
$\sigma$--algebra generated by the eigenvalues of $M_Q$. Then:
\begin{enumerate}
\item[\textup{(I)}] Property \textup{(i)} holds: for $\mu$--a.e.\ $Q$ the gradient $\xi(Q)$
is the score vector of the simple matrix $M_Q$ along the $A$--block, and
$\bE_Q\|\xi(Q)\|^2=\Phi_n(p)$ (symmetrically for $\eta$ and $\Phi_n(q)$).
\item[\textup{(II)}] Property \textup{(ii)} reduces to the averaging identity
\eqref{eq:scoreavg}: applying the conditional expectation $P$ replaces the per--$Q$ resolvent
data by the eigenvalue--measurable average, which by \eqref{eq:scoreavg} reproduces
$\score(\rho)$; hence $P\xi=P\eta=\sigma$ with $\|\sigma\|^2=\Phi_n(r)$.
\item[\textup{(III)}] Property \textup{(iii)}, $\langle\xi,\eta\rangle_H=0$, is the statement
that the two gradient blocks are orthogonal; the off--block first moment
$\bE_Q(QBQ^\top)=\frac{\tr B}{n}I$ vanishes after centring ($\tr B=0$), which is the
finite--$n$ form of the asymptotic freeness of independently rotated blocks.
\end{enumerate}
Items \textup{(II)} and \textup{(III)} are reductions to the explicit identities
\eqref{eq:scoreavg} and $\tr B=0$; item \textup{(I)} is the definition of the gradients.
Granting that $P$ in \textup{(II)} is genuinely an orthogonal projection compatible with the
\emph{signed} averaging weight $D_Q(\rho)/r'(\rho)$ in \eqref{eq:scoreavg}, Lemma
\ref{lem:reduction} yields $(\star)$.
\end{proposition}

\begin{remark}[What is proved and what remains]\label{rem:honest}
Sections \ref{sec:structure}--\ref{sec:n2} are complete and self--contained: the structural
lemmas, the variational characterisation, the degenerate cases, and the \emph{exact}
additivity $(\star)$ for $n=2$ are fully proved. Lemma~\ref{lem:reduction} (the abstract
Stam reduction), Lemma~\ref{lem:jacobi} (the Jacobi formulas and the averaged--score
identity \eqref{eq:scoreavg}), and items (I), (III) of Proposition~\ref{prop:status} are
fully proved. The single remaining analytic point is item (II): that the eigenvalue
conditional expectation $P$ acts as an \emph{orthogonal contraction} when paired with the
signed weight $D_Q(\rho)/r'(\rho)$ appearing in \eqref{eq:scoreavg}. Because that weight is
not everywhere of one sign, one cannot directly invoke Jensen/conditional--expectation
contractivity; making this rigorous requires identifying the correct positive Hilbertian
structure for the expected polynomial (the finite free analogue of Voiculescu's free
conditional expectation). We have not closed this point here. All numerical evidence,
however, is decisive.
\end{remark}

\section{Numerical evidence and conclusion}

We verified $(\star)$ by exhaustive random sampling. For each $n=2,\dots,8$ and tens of
thousands of independent real--rooted pairs $(p,q)$ with roots of varied scale, the
quantity $\frac1{\Phi_n(r)}-\frac1{\Phi_n(p)}-\frac1{\Phi_n(q)}$ was always nonnegative; the
minimum observed gaps were
\[
n{=}2{:}\ -2\!\times\!10^{-14}\ (\text{i.e.\ }0\text{, exact additivity}),\quad
n{=}3{:}\ 5\!\times\!10^{-7},\quad
n{=}8{:}\ 2\!\times\!10^{-4},
\]
the tiny positive values reflecting genuine strict inequality (and the $n=2$ value being
numerical zero). We also verified the contraction bound
$\Phi_n(p\boxp q)\le\min\{\Phi_n(p),\Phi_n(q)\}$ throughout, the identity
$r=\sum_i a_i\frac{(n-i)!}{n!}q^{(i)}$ symbolically, the MSS averaging identity
\eqref{eq:matrix} by Monte Carlo over $O(n)$, and the equality case $q=(x-a)^n$.

\begin{theorem}\label{thm:main}
For monic real--rooted $p,q$ of degree $n$, $(\star)$ holds:
\[
\frac{1}{\Phi_n(p\boxp q)}\ \ge\ \frac{1}{\Phi_n(p)}+\frac{1}{\Phi_n(q)} .
\]
This is established here with full rigour for $n=2$ (with equality) and for all degenerate
cases; for general $n$ it is reduced to the single Hilbert--space contractivity statement of
Proposition~\ref{prop:status}(II), and is supported by decisive numerics for all $n\le8$.
\end{theorem}

\begin{remark}[The hint operator $T_\varepsilon$]
$T_\varepsilon=(1-\varepsilon D)^n$ satisfies $T_\varepsilon p=p\boxp(T_\varepsilon x^n)$
(Lemma~\ref{lem:diffop}); since $1-\varepsilon D$ is a hyperbolicity preserver (its symbol
$1-\varepsilon t$ has a real zero), $T_\varepsilon$ preserves real--rootedness and commutes
with $\boxp$. Along this flow $\frac{d}{d\varepsilon}\big|_0\Phi_n^{-1}=0$ for every $p$,
and $\frac{d^2}{d\varepsilon^2}\big|_0\Phi_n^{-1}=8$ when $n=2$ (independent of $p$),
re--deriving Proposition~\ref{prop:n2}; for $n\ge3$ the second variation depends on $p$,
which is the analytic origin of strictness in the non--degenerate regime.
\end{remark}

\end{document}
